%------------------------------------------------------------------------------%
\documentclass[fleqn,utf8,aspectratio=169,12pt]{beamer}

\usepackage{gaal}

\title{Método de Gauss-Jordan}

%------------------------------------------------------------------------------%
\begin{document}

\addtitle

%------------------------------------------------------------------------------%
\section{Forma Escalonada Reduzida}

%------------------------------------------------------------------------------%
\begin{frame}{Forma Escalonada Reduzida}
  \onslide<+->{\emph{Forma escalonada reduzida}}
  \begin{itemize}[<+->]
    \item A matriz está em forma escalonada
    \item Cada pivô é igual a 1
    \item O pivô é o único elemento não nulo de sua coluna
  \end{itemize}

\bigskip
\onslide<+->{
Exemplo de matriz escalonada reduzida
\quad
\(
  \begin{bmatrix}
    \emph{1} & \z       & 2  & \z       \\
    \z       & \emph{1} & -1 & \z       \\
    \z       & \z       & \z & \emph{1}
  \end{bmatrix}
\)}
\end{frame}

%------------------------------------------------------------------------------%
\begin{frame}{Método de Gauss-Jordan}
  O processo para obter a Forma Escalonada Reduzida é chamado de \\[3mm]
  \emph{Eliminação de Gauss-Jordan}
\end{frame}

%------------------------------------------------------------------------------%
%------------------------------------------------------------------------------%
\begin{question}
Transforme a matriz
\[
  A =
  \begin{bmatrix}
    1 & 2 & 1 \\
    2 & 5 & 4 \\
    1 & 2 & 3
  \end{bmatrix}
\]
em sua forma escalonada reduzida
\end{question}

%------------------------------------------------------------------------------%
\begin{resolution}{Escalonamento}
\begin{columns}
\column{0.3\linewidth}
  \[\;\;\;
  \begin{bmatrix}
    1 & 2 & 1 \\
    2 & 5 & 4 \\
    1 & 2 & 3
  \end{bmatrix}
  \]
  \pause
\column{0.7\linewidth}
  Eliminamos os elementos abaixo do primeiro pivô
  \\ \pause
  \(
    L_2 \leftarrow L_2 - 2 L_1
  \)
  \\ \pause
  \(
    L_3 \leftarrow L_3 - L_1
  \)
\end{columns}
\vspace{-\baselineskip}
\[
  \pause
  L'_2
  \;=\;
  L_2 - 2 L_1
  \pause
  \;=\;
  \begin{bmatrix}
    2 & 5 & 4
  \end{bmatrix}
  - 2
  \begin{bmatrix}
    1 & 2 & 1
  \end{bmatrix}
  \pause
  \;=\;
  \begin{bmatrix}
    0 & 1 & 2
  \end{bmatrix}
\]
\[
  \pause
  L'_3
  \;=\;
  L_3 - L_1
  \pause
  \;=\;
  \begin{bmatrix}
    1 & 2 & 3
  \end{bmatrix}
  -
  \begin{bmatrix}
    1 & 2 & 1
  \end{bmatrix}
  \pause
  \;=\;
  \begin{bmatrix}
    0 & 0 & 2
  \end{bmatrix}
\]
\pause
\[
  \begin{bmatrix}
    1 & 2 & 1 \\
    0 & 1 & 2 \\
    0 & 0 & 2
  \end{bmatrix}
  \qquad\qquad
  \pause
  \text{Matriz escalonada}
\]
\end{resolution}

%------------------------------------------------------------------------------%
\begin{resolution}{Escalonamento}
\begin{columns}
\column{0.3\linewidth}
  \[\;\;\;
  \begin{bmatrix}
    1 & 2 & 1 \\
    0 & 1 & 2 \\
    0 & 0 & 2
  \end{bmatrix}
  \]
  \pause
\column{0.7\linewidth}
  Fazemos todos os pivôs iguais a 1
  \\ \pause
  \(
    L_3 \leftarrow \frac{1}{2} L_3
  \)
\end{columns}
\vspace{-\baselineskip}
\[
  \pause
  L'_3
  \;=\;
 \frac{1}{2} L_3
  \pause
  \;=\;
  \frac{1}{2}
  \begin{bmatrix}
    0 & 0 & 2
  \end{bmatrix}
  \pause
  \;=\;
  \begin{bmatrix}
    0 & 0 & 1
  \end{bmatrix}
\]
\pause
\[
  \begin{bmatrix}
    1 & 2 & 1 \\
    0 & 1 & 2 \\
    0 & 0 & 1
  \end{bmatrix}
\]
\end{resolution}

%------------------------------------------------------------------------------%
\begin{resolution}{Escalonamento}
\begin{columns}
\column{0.3\linewidth}
  \[\;\;\;
  \begin{bmatrix}
    1 & 2 & 1 \\
    0 & 1 & 2 \\
    0 & 0 & 1
  \end{bmatrix}
  \]
  \pause
\column{0.7\linewidth}
  Eliminamos os elementos acima do terceiro pivô
  \\ \pause
  \(
    L_2 \leftarrow L_2 - 2 L_3
  \)
  \\ \pause
  \(
    L_1 \leftarrow L_1 - L_3
  \)
\end{columns}
\vspace{-\baselineskip}
\[
  \pause
  L'_2
  \;=\;
  L_2 - 2 L_3
  \pause
  \;=\;
  \begin{bmatrix}
    0 & 1 & 2
  \end{bmatrix}
  - 2
  \begin{bmatrix}
    0 & 0 & 1
  \end{bmatrix}
  \pause
  \;=\;
  \begin{bmatrix}
    0 & 1 & 0
  \end{bmatrix}
\]
\[
  \pause
  L'_1
  \;=\;
  L_1 - L_3
  \pause
  \;=\;
  \begin{bmatrix}
    0 & 1 & 2
  \end{bmatrix}
  -
  \begin{bmatrix}
      0 & 0 & 1
  \end{bmatrix}
  \pause
  \;=\;
  \begin{bmatrix}
    1 & 2 & 0
  \end{bmatrix}
\]
\pause
\[
  \begin{bmatrix}
    1 & 2 & 0 \\
    0 & 1 & 0 \\
    0 & 0 & 1
  \end{bmatrix}
\]
\end{resolution}

%------------------------------------------------------------------------------%
\begin{resolution}{Escalonamento}
\begin{columns}
\column{0.3\linewidth}
  \[\;\;\;
  \begin{bmatrix}
    1 & 2 & 0 \\
    0 & 1 & 0 \\
    0 & 0 & 1
  \end{bmatrix}
  \]
  \pause
\column{0.7\linewidth}
  Eliminamos os elementos acima do segundo pivô
  \\ \pause
  \(
    L_1 \leftarrow L_1 - 2 L_2
  \)
\end{columns}
\vspace{-\baselineskip}
\[
  \pause
  L'_1
  \;=\;
  L_1 - 2 L_2
  \pause
  \;=\;
  \begin{bmatrix}
    1 & 2 & 0
  \end{bmatrix}
  \begin{bmatrix}
    0 & 1 & 0
  \end{bmatrix}
  \pause
  \;=\;
  \begin{bmatrix}
    1 & 0 & 0
  \end{bmatrix}
\]
\pause
\[
  \begin{bmatrix}
    1 & 0 & 0 \\
    0 & 1 & 0 \\
    0 & 0 & 1
  \end{bmatrix}
\]

\pause
Forma escalonada reduzida de $A$
\end{resolution}

%------------------------------------------------------------------------------%
\endinput

%------------------------------------------------------------------------------%
\begin{question}
Determine a forma escalonada reduzida da matriz
\[
  A=
  \begin{bmatrix}
    1  & 2 & -1 \\
    2  & 1 & 1  \\
    -1 & 3 & -2
  \end{bmatrix}
\]
\end{question}

%------------------------------------------------------------------------------%
\begin{resolution}{Escalonamento}
\begin{columns}
\column{0.3\linewidth}
  \[\;\;\;
  \begin{bmatrix}
    1  & 2 & -1 \\
    2  & 1 & 1  \\
    -1 & 3 & -2
  \end{bmatrix}
  \]
  \pause
\column{0.7\linewidth}
  Eliminamos os elementos abaixo do primeiro pivô
  \\ \pause
  \(
    L_2 \leftarrow L_2 -2 L_1
  \)
  \\ \pause
  \(
    L_3 \leftarrow L_3 + L_1
  \)
\end{columns}
\vspace{-\baselineskip}
\[
  \pause
  L'_2
  \;=\;
  L_2 -2 L_1
  \pause
  \;=\;
  \begin{bmatrix}
    2  & 1 & 1
  \end{bmatrix}
  -2
  \begin{bmatrix}
    1  & 2 & -1
  \end{bmatrix}
  \pause
  \;=\;
  \begin{bmatrix}
    0 & -3 & 3
  \end{bmatrix}
\]
\[
  \pause
  L'_3
  \;=\;
  L_3 + L_1
  \pause
  \;=\;
  \begin{bmatrix}
    -1 & 3 & -2
  \end{bmatrix}
  +
  \begin{bmatrix}
    1  & 2 & -1
  \end{bmatrix}
  \pause
  \;=\;
  \begin{bmatrix}
    0 & 5  & -3
  \end{bmatrix}
\]
\pause
\[
  \begin{bmatrix}
    1 & 2  & -1 \\
    0 & -3 & 3  \\
    0 & 5  & -3
  \end{bmatrix}
\]
\end{resolution}

%------------------------------------------------------------------------------%
\begin{resolution}{Escalonamento}
\begin{columns}
\column{0.3\linewidth}
  \[\;\;\;
  \begin{bmatrix}
    1 & 2  & -1 \\
    0 & -3 & 3  \\
    0 & 5  & -3
  \end{bmatrix}
  \]
  \pause
\column{0.7\linewidth}
  Fazendo com que o segundo pivô seja 1
  \\ \pause
  \(
  L_2 \leftarrow -\frac{1}{3}L_2
  \)
\end{columns}
\vspace{-\baselineskip}
\[
  \pause
  L'_2
  \;=\;
  \pause
  \;=\;
  -\frac{1}{3}L_2
  \begin{bmatrix}
    0 & -3 & 3
  \end{bmatrix}
  \pause
  \;=\;
  \begin{bmatrix}
    0 & 1 & -1
  \end{bmatrix}
\]
\pause
\[
  \begin{bmatrix}
    1 & 2 & -1 \\
    0 & 1 & -1 \\
    0 & 5 & -3
  \end{bmatrix}
\]
\end{resolution}

%------------------------------------------------------------------------------%
\begin{resolution}{Escalonamento}
\begin{columns}
\column{0.3\linewidth}
  \[\;\;\;
  \begin{bmatrix}
    1 & 2 & -1 \\
    0 & 1 & -1 \\
    0 & 5 & -3
  \end{bmatrix}
  \]
  \pause
\column{0.7\linewidth}
  Eliminamos os elementos abaixo do segundo pivô
  \\ \pause
  \(
    L_3 \leftarrow L_3 - 5 L_2
  \)
\end{columns}
\vspace{-\baselineskip}
\[
  \pause
  L'_3
  \;=\;
  L_3 - 5 L_2
  \pause
  \;=\;
  \begin{bmatrix}
    0 & 5 & -3
  \end{bmatrix}
  - 5
  \begin{bmatrix}
    0 & 1 & -1
  \end{bmatrix}
  \pause
  \;=\;
  \begin{bmatrix}
    0 & 0 & 2
  \end{bmatrix}
\]
\pause
\[
  \begin{bmatrix}
    1 & 2 & -1 \\
    0 & 1 & -1 \\
    0 & 0 & 2
  \end{bmatrix}
\]
\end{resolution}

%------------------------------------------------------------------------------%
\begin{resolution}{Escalonamento}
\begin{columns}
\column{0.3\linewidth}
  \[\;\;\;
  \begin{bmatrix}
    1 & 2 & -1 \\
    0 & 1 & -1 \\
    0 & 0 & 2
  \end{bmatrix}
  \]
  \pause
\column{0.7\linewidth}
  Fazendo com que o terceiro pivô seja 1
  \\ \pause
  \(
    L_3 \leftarrow \frac{1}{2} L_3
  \)
\end{columns}
\vspace{-\baselineskip}
\[
  \pause
  L'_3
  \;=\;
  \frac{1}{2} L_3
  \pause
  \;=\;
  \frac{1}{2}
  \begin{bmatrix}
    0 & 0 & 2
  \end{bmatrix}
  \pause
  \;=\;
  \begin{bmatrix}
    0 & 0 & 1
  \end{bmatrix}
\]
\pause
\[
  \begin{bmatrix}
    1 & 2 & -1 \\
    0 & 1 & -1 \\
    0 & 0 & 1
  \end{bmatrix}
\]
\end{resolution}

%------------------------------------------------------------------------------%
\begin{resolution}{Escalonamento}
\begin{columns}
\column{0.3\linewidth}
  \[\;\;\;
  \begin{bmatrix}
    1 & 2 & -1 \\
    0 & 1 & -1 \\
    0 & 0 & 1
  \end{bmatrix}
  \]
  \pause
\column{0.7\linewidth}
  Eliminamos os elementos acima do terceiro pivô
  \\ \pause
  \(
    L_2 \leftarrow L_2 + L_3
  \)
  \\ \pause
  \(
    L_1 \leftarrow L_1 + L_3
  \)
\end{columns}
\vspace{-\baselineskip}
\[
  \pause
  L'_2
  \;=\;
  L_2 + L_3
  \pause
  \;=\;
  \begin{bmatrix}
    0 & 1 & -1
  \end{bmatrix}
  +
  \begin{bmatrix}
    0 & 0 & 1
  \end{bmatrix}
  \pause
  \;=\;
  \begin{bmatrix}
    0 & 1 & 0
  \end{bmatrix}
\]
\[
  \pause
  L'_1
  \;=\;
  L_1 + L_3
  \pause
  \;=\;
  \begin{bmatrix}
    1 & 2 & -1
  \end{bmatrix}
  +
  \begin{bmatrix}
    0 & 0 & 1
  \end{bmatrix}
  \pause
  \;=\;
  \begin{bmatrix}
    1 & 2 & 0
  \end{bmatrix}
\]
\pause
\[
  \begin{bmatrix}
    1 & 2 & 0 \\
    0 & 1 & 0 \\
    0 & 0 & 1
  \end{bmatrix}
\]
\end{resolution}

%------------------------------------------------------------------------------%
\begin{resolution}{Escalonamento}
\begin{columns}
\column{0.3\linewidth}
  \[\;\;\;
  \begin{bmatrix}
    1 & 2 & 0 \\
    0 & 1 & 0 \\
    0 & 0 & 1
  \end{bmatrix}
  \]
  \pause
\column{0.7\linewidth}
  Eliminamos os elementos acima do segundo pivô
  \\ \pause
  \(
    L_1 \leftarrow L_1 - 2 L_2
  \)
\end{columns}
\vspace{-\baselineskip}
\[
  \pause
  L'_1
  \;=\;
  L_1 - 2 L_2
  \pause
  \;=\;
  \begin{bmatrix}
    1 & 2 & 0
  \end{bmatrix}
  - 2
  \begin{bmatrix}
    0 & 1 & 0
  \end{bmatrix}
  \pause
  \;=\;
  \begin{bmatrix}
    1 & 0 & 0
  \end{bmatrix}
\]
\pause
\[
  \begin{bmatrix}
    1 & 0 & 0 \\
    0 & 1 & 0 \\
    0 & 0 & 1
  \end{bmatrix}
\]

\pause
Forma escalonada reduzida da matriz $A$
\end{resolution}

%------------------------------------------------------------------------------%
\endinput

%------------------------------------------------------------------------------%
\begin{question}
Resolva o sistema linear pelo método de Gauss-Jordan
\[
  \systeme[sort={x,y,z,w}]{
     x + 2y - z + w =  4,
    2x + 5y     + w =  7,
    -x -  y +2z - w = -1,
    3x + 7y -z  + w = 10
  }
\]
\end{question}

%------------------------------------------------------------------------------%
\begin{resolution}{Matriz Aumentada}
  Matriz aumentada
  \[
    \left[\begin{array}{rrrr|r}
      1  & 2  & -1 & 1  & 4  \\
      2  & 5  & 0  & 1  & 7  \\
      -1 & -1 & 2  & -1 & -1 \\
      3  & 7  & -1 & 1  & 10
    \end{array}\right]
  \]
\end{resolution}

%------------------------------------------------------------------------------%
\begin{resolution}{Escalonamento do Sistema}
\begin{columns}
\column{0.4\linewidth}
  \[\;\;\,
    \left[\begin{array}{rrrr|r}
      1  & 2  & -1 & 1  & 4  \\
      2  & 5  & 0  & 1  & 7  \\
      -1 & -1 & 2  & -1 & -1 \\
      3  & 7  & -1 & 1  & 10
    \end{array}\right]
  \]
  \pause
\column{0.6\linewidth}
  Eliminar os elementos abaixo do primeiro pivô
  \\ \pause
  \(
    L_2 \leftarrow L_2 - 2 L_1
  \)
  \\ \pause
  \(
    L_3 \leftarrow L_3 + L_1
  \)
  \\ \pause
  \(
    L_4 \leftarrow L_4 - 3 L_1
  \)
\end{columns}
\vspace{-\baselineskip}
\[
  \pause
  \hspace{-3.2em}
  L'_2
  \;=\;
  L_2 - 2 L_1
  \pause
  \;=\;
  \left[\begin{array}{rrrr|r}
    2  & 5  & 0  & 1  & 7
  \end{array}\right]
  -2
  \left[\begin{array}{rrrr|r}
    1  & 2  & -1 & 1  & 4
  \end{array}\right]
  \pause
  \;=\;
  \left[\begin{array}{rrrr|r}
    0 & 1 & 2  & -1 & -1
  \end{array}\right]
\]
\[
  \pause
  \hspace{-3.2em}
  L'_3
  \;=\;
  L_3 + L_1
  \pause
  \;=\;
  \left[\begin{array}{rrrr|r}
    -1 & -1 & 2  & -1 & -1
  \end{array}\right]
  +
  \left[\begin{array}{rrrr|r}
    1  & 2  & -1 & 1  & 4
  \end{array}\right]
  \pause
  \;=\;
  \left[\begin{array}{rrrr|r}
    0 & 1 & 1  & 0  & 3
  \end{array}\right]
\]
\[
  \pause
  \hspace{-3.2em}
  L'_4
  \;=\;
  L_4 - 3 L_1
  \pause
  \;=\;
  \left[\begin{array}{rrrr|r}
    3  & 7  & -1 & 1  & 10
  \end{array}\right]
  -3
  \left[\begin{array}{rrrr|r}
    1  & 2  & -1 & 1  & 4
  \end{array}\right]
  \pause
  \;=\;
  \left[\begin{array}{rrrr|r}
    0 & 1 & 2  & -2 & -2
  \end{array}\right]
\]
\[
  \pause
  \left[\begin{array}{rrrr|r}
    1 & 2 & -1 & 1  & 4  \\
    0 & 1 & 2  & -1 & -1 \\
    0 & 1 & 1  & 0  & 3  \\
    0 & 1 & 2  & -2 & -2
  \end{array}\right]
\]
\end{resolution}

%------------------------------------------------------------------------------%
\begin{resolution}{Escalonamento do Sistema}
\begin{columns}
\column{0.4\linewidth}
  \[\;\;\,
    \left[\begin{array}{rrrr|r}
      1 & 2 & -1 & 1  & 4  \\
      0 & 1 & 2  & -1 & -1 \\
      0 & 1 & 1  & 0  & 3  \\
      0 & 1 & 2  & -2 & -2
    \end{array}\right]
  \]
  \pause
\column{0.6\linewidth}
  Eliminar os elementos abaixo do segundo pivô
  \\ \pause
  \(
    L_3 \leftarrow L_3 - L_2
  \)
  \\ \pause
  \(
    L_4 \leftarrow L_4 - L_2
  \)
\end{columns}
\vspace{-\baselineskip}
\[
  \pause
  \hspace{-3.2em}
  L'_3
  \;=\;
  L_3 - L_2
  \pause
  \;=\;
  \left[\begin{array}{rrrr|r}
    0 & 1 & 1  & 0  & 3
  \end{array}\right]
  -
  \left[\begin{array}{rrrr|r}
    0 & 1 & 2  & -1 & -1
  \end{array}\right]
  \pause
  \;=\;
  \left[\begin{array}{rrrr|r}
    0 & 0 & -1 & 1  & 4
  \end{array}\right]
\]
\[
  \pause
  \hspace{-3.2em}
  L'_4
  \;=\;
  L_4 - L_2
  \pause
  \;=\;
  \left[\begin{array}{rrrr|r}
    0 & 1 & 2  & -2 & -2
  \end{array}\right]
  -
  \left[\begin{array}{rrrr|r}
    0 & 1 & 2  & -1 & -1
  \end{array}\right]
  \pause
  \;=\;
  \left[\begin{array}{rrrr|r}
    0 & 0 & 0  & -1 & -1
  \end{array}\right]
\]
\[
  \pause
  \left[\begin{array}{rrrr|r}
    1 & 2 & -1 & 1  & 4  \\
    0 & 1 & 2  & -1 & -1 \\
    0 & 0 & -1 & 1  & 4  \\
    0 & 0 & 0  & -1 & -1
  \end{array}\right]
\]
\end{resolution}

%------------------------------------------------------------------------------%
\begin{resolution}{Escalonamento do Sistema}
\begin{columns}
\column{0.4\linewidth}
  \[\;\;\,
    \left[\begin{array}{rrrr|r}
      1 & 2 & -1 & 1  & 4  \\
      0 & 1 & 2  & -1 & -1 \\
      0 & 0 & -1 & 1  & 4  \\
      0 & 0 & 0  & -1 & -1
    \end{array}\right]
  \]
  \pause
\column{0.6\linewidth}
  Transformar os pivôs em 1
  \\ \pause
  \(
    L_3 \leftarrow - L_3
  \)
  \\ \pause
  \(
    L_4 \leftarrow - L_4
  \)
\end{columns}
\vspace{-\baselineskip}
\[
  \pause
  L'_3
  \;=\;
  - L_3
  \pause
  \;=\;
  -
  \left[\begin{array}{rrrr|r}
    0 & 0 & -1 & 1  & 4
  \end{array}\right]
  \pause
  \;=\;
  \left[\begin{array}{rrrr|r}
    0 & 0 & 1  & -1 & -4
  \end{array}\right]
\]
\[
  \pause
  L'_4
  \;=\;
  - L_4
  \pause
  \;=\;
  -
  \left[\begin{array}{rrrr|r}
    0 & 0 & 0  & -1 & -1
  \end{array}\right]
  \pause
  \;=\;
  \left[\begin{array}{rrrr|r}
    0 & 0 & 0  & 1  & 1
  \end{array}\right]
\]
\[
  \pause
  \left[\begin{array}{rrrr|r}
    1 & 2 & -1 & 1  & 4  \\
    0 & 1 & 2  & -1 & -1 \\
    0 & 0 & 1  & -1 & -4 \\
    0 & 0 & 0  & 1  & 1
  \end{array}\right]
\]
\end{resolution}

%------------------------------------------------------------------------------%
\begin{resolution}{Escalonamento do Sistema}
\begin{columns}
\column{0.4\linewidth}
  \[\;\;\,
    \left[\begin{array}{rrrr|r}
      1 & 2 & -1 & 1  & 4  \\
      0 & 1 & 2  & -1 & -1 \\
      0 & 0 & 1  & -1 & -4 \\
      0 & 0 & 0  & 1  & 1
    \end{array}\right]
  \]
  \pause
\column{0.6\linewidth}
  Eliminar os elementos acima do quarto pivô
  \\ \pause
  \(
    L_3 \leftarrow L_3 + L_4
  \)
  \\ \pause
  \(
    L_2 \leftarrow L_2 + L_4
  \)
  \\ \pause
  \(
    L_1 \leftarrow L_1 - L_4
  \)
\end{columns}
\vspace{-\baselineskip}
\[
  \pause
  \hspace{-1.2em}
  L'_3
  \;=\;
  L_3 + L_4
  \pause
  \;=\;
  \left[\begin{array}{rrrr|r}
    0 & 0 & 1  & -1 & -4
  \end{array}\right]
  +
  \left[\begin{array}{rrrr|r}
    0 & 0 & 0  & 1  & 1
  \end{array}\right]
  \pause
  \;=\;
  \left[\begin{array}{rrrr|r}
    0 & 0 & 1  & 0 & -3
  \end{array}\right]
\]
\[
  \pause
  \hspace{-1.2em}
  L'_2
  \;=\;
  L_2 + L_4
  \pause
  \;=\;
  \left[\begin{array}{rrrr|r}
    0 & 1 & 2  & -1 & -1
  \end{array}\right]
  +
  \left[\begin{array}{rrrr|r}
    0 & 0 & 0  & 1  & 1
  \end{array}\right]
  \pause
  \;=\;
  \left[\begin{array}{rrrr|r}
    0 & 1 & 2  & 0 & 0
  \end{array}\right]
\]
\[
  \pause
  \hspace{-1.2em}
  L'_1
  \;=\;
  L_1 - L_4
  \pause
  \;=\;
  \left[\begin{array}{rrrr|r}
    1 & 2 & -1 & 1  & 4
  \end{array}\right]
  -
  \left[\begin{array}{rrrr|r}
    0 & 0 & 0  & 1  & 1
  \end{array}\right]
  \pause
  \;=\;
  \left[\begin{array}{rrrr|r}
    1 & 2 & -1 & 0 & 3
  \end{array}\right]
\]
\[
  \pause
  \left[\begin{array}{rrrr|r}
    1 & 2 & -1 & 0 & 3  \\
    0 & 1 & 2  & 0 & 0  \\
    0 & 0 & 1  & 0 & -3 \\
    0 & 0 & 0  & 1 & 1
  \end{array}\right]
\]
\end{resolution}

%------------------------------------------------------------------------------%
\begin{resolution}{Escalonamento do Sistema}
\begin{columns}
\column{0.4\linewidth}
  \[\;\;\,
    \left[\begin{array}{rrrr|r}
      1 & 2 & -1 & 0 & 3  \\
      0 & 1 & 2  & 0 & 0  \\
      0 & 0 & 1  & 0 & -3 \\
      0 & 0 & 0  & 1 & 1
    \end{array}\right]
  \]
  \pause
\column{0.6\linewidth}
  Eliminar os elementos acima do terceiro pivô
  \\ \pause
  \(
    L_2 \leftarrow L_2 - 2 L_3
  \)
  \\ \pause
  \(
    L_1 \leftarrow L_1 + L_3
  \)
\end{columns}
\vspace{-\baselineskip}
\[
  \pause
  \hspace{-1.2em}
  L'_2
  \;=\;
  L_2 - 2 L_3
  \pause
  \;=\;
  \left[\begin{array}{rrrr|r}
    0 & 1 & 2  & 0 & 0
  \end{array}\right]
  - 2
  \left[\begin{array}{rrrr|r}
    0 & 0 & 1  & 0 & -3
  \end{array}\right]
  \pause
  \;=\;
  \left[\begin{array}{rrrr|r}
    0 & 1 & 0 & 0 & 6
  \end{array}\right]
\]
\[
  \pause
  \hspace{-1.2em}
  L'_1
  \;=\;
  L_1 + L_3
  \pause
  \;=\;
  \left[\begin{array}{rrrr|r}
    1 & 2 & -1 & 0 & 3
  \end{array}\right]
  +
  \left[\begin{array}{rrrr|r}
    0 & 0 & 1  & 0 & -3
  \end{array}\right]
  \pause
  \;=\;
  \left[\begin{array}{rrrr|r}
    1 & 2 & 0 & 0 & 0
  \end{array}\right]
\]
\[
  \pause
  \left[\begin{array}{rrrr|r}
    1 & 2 & 0 & 0 &  0 \\
    0 & 1 & 0 & 0 &  6 \\
    0 & 0 & 1 & 0 & -3 \\
    0 & 0 & 0 & 1 &  1
  \end{array}\right]
\]
\end{resolution}

%------------------------------------------------------------------------------%
\begin{resolution}{Escalonamento do Sistema}
\begin{columns}
\column{0.4\linewidth}
  \[\;\;\,
    \left[\begin{array}{rrrr|r}
      1 & 2 & 0 & 0 &  0 \\
      0 & 1 & 0 & 0 &  6 \\
      0 & 0 & 1 & 0 & -3 \\
      0 & 0 & 0 & 1 &  1
    \end{array}\right]
  \]
  \pause
\column{0.6\linewidth}
  Eliminar os elementos acima do segundo pivô
  \\ \pause
  \(
    L_1 \leftarrow L_1 - 2 L_2
  \)
\end{columns}
\vspace{-\baselineskip}
\[
  \pause
  \hspace{-2em}
  L'_1
  \;=\;
  L_1 - 2 L_2
  \pause
  \;=\;
  \left[\begin{array}{rrrr|r}
    1 & 2 & 0 & 0 &  0
  \end{array}\right]
  - 2
  \left[\begin{array}{rrrr|r}
    0 & 1 & 0 & 0 &  6
  \end{array}\right]
  \pause
  \;=\;
  \left[\begin{array}{rrrr|r}
    1 & 0 & 0 & 0 &  -12
  \end{array}\right]
\]
\[
  \pause
  \left[\begin{array}{rrrr|r}
    1 & 0 & 0 & 0 & -12 \\
    0 & 1 & 0 & 0 &  6 \\
    0 & 0 & 1 & 0 & -3 \\
    0 & 0 & 0 & 1 & 1
  \end{array}\right]
\]
\end{resolution}

%------------------------------------------------------------------------------%
\begin{resolution}{Solução}
A solução do sistema é
\(
  {\bm x}
  \;=\;
  \begin{bmatrix}
    x \\ y \\ z \\ w
  \end{bmatrix}
  \;=\;
  \begin{bmatrix}
    -12 \\ 6 \\ -3 \\ 1
  \end{bmatrix}
\)
\end{resolution}

%------------------------------------------------------------------------------%
\endinput



%------------------------------------------------------------------------------%
\begin{question}
  Resolva o sistema linear pelo método de Gauss-Jordan
%  \[
%    \systeme[sort={x,y,z,w,v}]{
%       x + 2y -  z +  w +  v =  3,
%      2x + 5y +  z + 2w +  v =  8,
%      -x -  y + 2z -  w + 2v = -1,
%            y +  z      + 3v =  2
%    }
%  \]
  \[
    \left\{\begin{array}{rrrrrr}
       x & + 2y & -  z & +  w & +  v & = \hpm 3 \\
      2x & + 5y & +  z & + 2w & +  v & = \hpm 8 \\
      -x & -  y & + 2z & -  w & + 2v & =     -1 \\
         &    y & +  z &      & + 3v & = \hpm 2
    \end{array}\right.
  \]
\end{question}

%------------------------------------------------------------------------------%
\begin{resolution}{Matriz Aumentada}
  Matriz aumentada
  \[
    \begin{amatrix}{5}
      1  & 2  & -1 & 1  & 1 & 3  \\
      2  & 5  & 1  & 2  & 1 & 8  \\
      -1 & -1 & 2  & -1 & 2 & -1 \\
      0  & 1  & 1  & 0  & 3 & 2
    \end{amatrix}
  \]
\end{resolution}

%------------------------------------------------------------------------------%
\begin{resolution}{Escalonamento do Sistema}
\begin{columns}
\column{0.45\linewidth}
  \[
    \begin{amatrix}{5}
      1  & 2  & -1 & 1  & 1 & 3  \\
      2  & 5  & 1  & 2  & 1 & 8  \\
      -1 & -1 & 2  & -1 & 2 & -1 \\
      0  & 1  & 1  & 0  & 3 & 2
    \end{amatrix}
  \]
  \pause
\column{0.58\linewidth}
  Eliminar os elementos abaixo do primeiro pivô
  \\ \pause
  \(
    L_2 \leftarrow L_2 - 2 L_1
  \)
  \\ \pause
  \(
    L_3 \leftarrow L_3 + L_1
  \)
\end{columns}
\vspace{-\baselineskip}
\[
  \pause
  \hspace{-2.5em}
  L'_2
  \pause
  \;=\;
  \begin{amatrix}{5}
    2  & 5  & 1  & 2  & 1 & 8
  \end{amatrix}
  -2
  \begin{amatrix}{5}
    1  & 2  & -1 & 1  & 1 & 3
  \end{amatrix}
  \pause
  \;=\;
  \begin{amatrix}{5}
    0 & 1 & 3  & 0 & -1 & 2
  \end{amatrix}
\]
\[
  \pause
  \hspace{-2.5em}
  L'_3
  \pause
  \;=\;
  \begin{amatrix}{5}
    -1 & -1 & 2  & -1 & 2 & -1
  \end{amatrix}
  +
  \begin{amatrix}{5}
    1  & 2  & -1 & 1  & 1 & 3
  \end{amatrix}
  \pause
  \;=\;
  \begin{amatrix}{5}
    0 & 1 & 1  & 0 & 3  & 2
  \end{amatrix}
\]
\pause
\[
  \begin{amatrix}{5}
    1 & 2 & -1 & 1 & 1  & 3 \\
    0 & 1 & 3  & 0 & -1 & 2 \\
    0 & 1 & 1  & 0 & 3  & 2 \\
    0 & 1 & 1  & 0 & 3  & 2
  \end{amatrix}
\]
\end{resolution}

%------------------------------------------------------------------------------%
\begin{resolution}{Escalonamento do Sistema}
\begin{columns}
\column{0.45\linewidth}
  \[
    \begin{amatrix}{5}
      1 & 2 & -1 & 1 & 1  & 3 \\
      0 & 1 & 3  & 0 & -1 & 2 \\
      0 & 1 & 1  & 0 & 3  & 2 \\
      0 & 1 & 1  & 0 & 3  & 2
    \end{amatrix}
  \]
  \pause
\column{0.58\linewidth}
  Eliminar os elementos abaixo do segundo pivô
  \\ \pause
  \(
    L_3 \leftarrow L_3 - L_2
  \)
  \\ \pause
  \(
    L_4 \leftarrow L_4 - L_2
  \)
\end{columns}
\vspace{-\baselineskip}
\[
  \pause
  \hspace{-1.5em}
  L'_3
  \pause
  \;=\;
  \begin{amatrix}{5}
    0 & 1 & 1  & 0 & 3  & 2
  \end{amatrix}
  -
  \begin{amatrix}{5}
    0 & 1 & 3  & 0 & -1 & 2
  \end{amatrix}
  \pause
  \;=\;
  \begin{amatrix}{5}
    0 & 0 & -2 & 0 & 4  & 0
  \end{amatrix}
\]
\[
  \pause
  \hspace{-1.5em}
  L'_4
  \pause
  \;=\;
  \begin{amatrix}{5}
    0 & 1 & 1  & 0 & 3  & 2
  \end{amatrix}
  -
  \begin{amatrix}{5}
    0 & 1 & 3  & 0 & -1 & 2
  \end{amatrix}
  \pause
  \;=\;
  \begin{amatrix}{5}
    0 & 0 & -2 & 0 & 4  & 0
  \end{amatrix}
\]
\pause
\[
  \begin{amatrix}{5}
    1 & 2 & -1 & 1 & 1  & 3 \\
    0 & 1 & 3  & 0 & -1 & 2 \\
    0 & 0 & -2 & 0 & 4  & 0 \\
    0 & 0 & -2 & 0 & 4  & 0
  \end{amatrix}
\]
\end{resolution}

%------------------------------------------------------------------------------%
\begin{resolution}{Escalonamento do Sistema}
\begin{columns}
\column{0.45\linewidth}
  \[
    \begin{amatrix}{5}
      1 & 2 & -1 & 1 & 1  & 3 \\
      0 & 1 & 3  & 0 & -1 & 2 \\
      0 & 0 & -2 & 0 & 4  & 0 \\
      0 & 0 & -2 & 0 & 4  & 0
    \end{amatrix}
  \]
  \pause
\column{0.58\linewidth}
  Tornar o terceiro pivô igual a 1
  \\ \pause
  \(
    L_3 \leftarrow -\frac{1}{2}L_3
  \)
\end{columns}
\vspace{-\baselineskip}
\[
  \pause
  \hspace{-1.5em}
  L'_3
  \pause
  \;=\;
  -\frac{1}{2}
  \begin{amatrix}{5}
    0 & 0 & -2 & 0 & 4  & 0
  \end{amatrix}
  \pause
  \;=\;
  \begin{amatrix}{5}
    0 & 0 & 1  & 0 & -2 & 0
  \end{amatrix}
\]
\pause
\[
  \begin{amatrix}{5}
    1 & 2 & -1 & 1 & 1  & 3 \\
    0 & 1 & 3  & 0 & -1 & 2 \\
    0 & 0 & 1  & 0 & -2 & 0 \\
    0 & 0 & -2 & 0 & 4  & 0
  \end{amatrix}
\]
\end{resolution}

%------------------------------------------------------------------------------%
\begin{resolution}{Escalonamento do Sistema}
\begin{columns}
\column{0.4\linewidth}
  \[
    \begin{amatrix}{5}
      1 & 2 & -1 & 1 & 1  & 3 \\
      0 & 1 & 3  & 0 & -1 & 2 \\
      0 & 0 & 1  & 0 & -2 & 0 \\
      0 & 0 & -2 & 0 & 4  & 0
    \end{amatrix}
  \]
  \pause
\column{0.6\linewidth}
  Eliminar os elementos abaixo do terceiro pivô
  \\ \pause
  \(
    L_4 \leftarrow L_4 +2 L_3
  \)
\end{columns}
\vspace{-\baselineskip}
\[
  \pause
  \hspace{-1.5em}
  L'_4
  \pause
  \;=\;
  \begin{amatrix}{5}
    0 & 0 & -2 & 0 & 4  & 0
  \end{amatrix}
  +2
  \begin{amatrix}{5}
    0 & 0 & 1  & 0 & -2 & 0
  \end{amatrix}
  \pause
  \;=\;
  \begin{amatrix}{5}
    0 & 0 & 0 & 0 & 0 & 0
  \end{amatrix}
\]
\pause
\[
  \begin{amatrix}{5}
    1 & 2 & -1 & 1 & 1  & 3 \\
    0 & 1 & 3  & 0 & -1 & 2 \\
    0 & 0 & 1  & 0 & -2 & 0 \\
    0 & 0 & 0  & 0 &  0 & 0
  \end{amatrix}
\]
\end{resolution}

%------------------------------------------------------------------------------%
\begin{resolution}{Escalonamento do Sistema}
\begin{columns}
\column{0.4\linewidth}
  \[
    \begin{amatrix}{5}
      1 & 2 & -1 & 1 & 1  & 3 \\
      0 & 1 & 3  & 0 & -1 & 2 \\
      0 & 0 & 1  & 0 & -2 & 0 \\
      0 & 0 & 0  & 0 &  0 & 0
    \end{amatrix}
  \]
  \pause
\column{0.6\linewidth}
  Eliminar os elementos acima do terceiro pivô
  \\ \pause
  \(
    L_2 \leftarrow L_2 -3 L_3
  \)
  \\ \pause
  \(
    L_1 \leftarrow L_1 + L_3
  \)
\end{columns}
\vspace{-\baselineskip}
\[
  \pause
  \hspace{-1.5em}
  L'_2
  \pause
  \;=\;
  \begin{amatrix}{5}
    0 & 1 & 3  & 0 & -1 & 2
  \end{amatrix}
  -3
  \begin{amatrix}{5}
    0 & 0 & 1  & 0 & -2 & 0
  \end{amatrix}
  \pause
  \;=\;
  \begin{amatrix}{5}
    0 & 1 & 0 & 0 & 5  & 2
  \end{amatrix}
\]
\[
  \pause
  \hspace{-1.5em}
  L'_1
  \pause
  \;=\;
  \begin{amatrix}{5}
    1 & 2 & -1 & 1 & 1  & 3
  \end{amatrix}
  +
  \begin{amatrix}{5}
    0 & 0 & 1  & 0 & -2 & 0
  \end{amatrix}
  \pause
  \;=\;
  \begin{amatrix}{5}
    1 & 2 & 0 & 1 & -1 & 3
  \end{amatrix}
\]
\pause
\[
  \begin{amatrix}{5}
    1 & 2 & 0 & 1 & -1 & 3 \\
    0 & 1 & 0 & 0 & 5  & 2 \\
    0 & 0 & 1 & 0 & -2 & 0 \\
    0 & 0 & 0 & 0 & 0  & 0
  \end{amatrix}
\]
\end{resolution}

%------------------------------------------------------------------------------%
\begin{resolution}{Escalonamento do Sistema}
\begin{columns}
\column{0.4\linewidth}
  \[
    \begin{amatrix}{5}
      1 & 2 & 0 & 1 & -1 & 3 \\
      0 & 1 & 0 & 0 & 5  & 2 \\
      0 & 0 & 1 & 0 & -2 & 0 \\
      0 & 0 & 0 & 0 & 0  & 0
    \end{amatrix}
  \]
  \pause
\column{0.6\linewidth}
  Eliminar os elementos acima do segundo pivô
  \\ \pause
  \(
    L_1 \leftarrow L_1 - 2 L_2
  \)
\end{columns}
\vspace{-\baselineskip}
\[
  \pause
  \hspace{-1.5em}
  L'_1
  \pause
  \;=\;
  \begin{amatrix}{5}
    1 & 2 & 0 & 1 & -1 & 3
  \end{amatrix}
  -2
  \begin{amatrix}{5}
    0 & 1 & 0 & 0 & 5  & 2
  \end{amatrix}
  \pause
  \;=\;
  \begin{amatrix}{5}
    1 & 0 & 0 & 1 & -11 & -1
  \end{amatrix}
\]
\pause
\[
  \begin{amatrix}{5}
    1 & 0 & 0 & 1 & -11 & -1 \\
    0 & 1 & 0 & 0 & 5   & 2  \\
    0 & 0 & 1 & 0 & -2  & 0  \\
    0 & 0 & 0 & 0 & 0   & 0
  \end{amatrix}
\]
\end{resolution}

%------------------------------------------------------------------------------%
\begin{resolution}{Sistema Linear Reduzido}
%  \[
%    \systeme[sort={x,y,z,w,v}]{
%      x +  w -11v = -1,
%      y + 5v      = 2,
%      z - 2v      = 0
%    }
%  \]
  \[
    \left\{\begin{array}{llllll}
      x &   &   & + w & - 11v & = -1 \\
        & y &   &     & +  5v & = \hpm 2 \\
        &   & z &     & -  2v & = \hpm 0
    \end{array}\right.
  \]
  \pause
  As incógnitas $w$ e $v$ não correspondem a pivôs

  \pause
  portanto são variáveis livres
  \begin{align*}
    w & = s \\
    v & = t
  \end{align*}
  para quaisquer valore reais $t$ e $s$
\end{resolution}

%------------------------------------------------------------------------------%
\begin{resolution}{Solução}
  \[
    \left\{\begin{array}{llllll}
      x &   &   & + w & - 11v & = -1 \\
        & y &   &     & +  5v & = \hpm 2 \\
        &   & z &     & -  2v & = \hpm 0
    \end{array}\right.
  \]
%  \[
%    \systeme[sort={x,y,z,w,v}]{
%      x +  w -11v = -1,
%      y + 5v      = 2,
%      z - 2v      = 0
%    }
%  \]
  \pause
  \vspace{-\baselineskip}
  \begin{columns}[t]
  \column{0.35\linewidth}
    \begin{align*}
      x & = 11v - w -1 \\[-1.5ex]
      y & = 2 - 5v     \\[-1.5ex]
      z & = 2v
    \end{align*}
    \pause
  \column{0.6\linewidth}
    \begin{align*}
      x & = 11t - s - 1 \\[-1.5ex]
      y & = 2 - 5t      \\[-1.5ex]
      z & = 2t          \qquad\qquad\qquad \forall (s, t) \in \R^2 \\[-1.5ex]
      w & = s           \\[-1.5ex]
      v & = t
    \end{align*}
  \end{columns}
\end{resolution}

%------------------------------------------------------------------------------%
\endinput


%------------------------------------------------------------------------------%
%\section{Lista Mínima}
%
%%------------------------------------------------------------------------------%
%\begin{frame}{Lista Mínima}
%  \begin{enumerate}
%    \item
%  \end{enumerate}
%
%  \vfill
%  Atenção: A prova é baseada no livro, não nas apresentações
%\end{frame}

%------------------------------------------------------------------------------%
\end{document}
%------------------------------------------------------------------------------%
